\documentclass[11pt]{article}
\usepackage[utf8]{inputenc}
\usepackage[T1]{fontenc}
\usepackage{amsmath,amssymb}
\usepackage{graphicx}
\usepackage{booktabs}
\usepackage{hyperref}
\usepackage[margin=1in]{geometry}
\usepackage{natbib}
\bibliographystyle{unsrtnat}

\title{Stable Population Structure in Europe Since the Iron Age, Despite High Mobility}
\author{Authors}
\date{}

\begin{document}
\maketitle

\begin{abstract}
Ancient DNA research in the past decade has revealed that European population structure changed dramatically in the prehistoric period (14,000--3,000 years before present, YBP), reflecting the widespread introduction of Neolithic farmer and Bronze Age Steppe ancestries. However, little is known about how population structure changed from the historical period onward (3,000 YBP -- present). To address this, we collected whole genomes from 204 individuals from Europe and the Mediterranean, many of which are the first historical period genomes from their region (e.g., Armenia and France). We found that most regions show remarkable inter-individual heterogeneity. At least 7\% of historical individuals carry ancestry uncommon in the region where they were sampled, some indicating cross-Mediterranean contacts. Despite this high level of mobility, overall population structure across western Eurasia is relatively stable through the historical period up to the present, mirroring geography. We show that, under standard population genetics models with local panmixia, the observed level of dispersal would lead to a collapse of population structure. Persistent population structure thus suggests a lower effective migration rate than indicated by the observed dispersal. We hypothesize that this phenomenon can be explained by extensive transient dispersal arising from drastically improved transportation networks and the Roman Empire's mobilization of people for trade, labor, and military.
\end{abstract}

\section{Introduction}

Ancient DNA (aDNA) sequencing has provided immense insight into previously unanswered questions about human population history \citep{Allentoft2015,Haak2015,Lazaridis2014}. Initially, sequencing efforts focused on identifying the main ancestry groups and transitions during prehistoric times, for which there is no written record. Recently, aDNA sampling has expanded to more recent times, allowing the study of movements of people using genetic data alongside the well-studied historical record \citep{Antonio2019,Clemente2021}.

Prehistoric ancient DNA studies revealed that European populations experienced two major waves of migration: the spread of Neolithic farming from Anatolia beginning around 9,000 years ago, and the influx of Steppe-related ancestry during the Bronze Age \citep{Haak2015,Lazaridis2016,Mathieson2015}. These migrations fundamentally reshaped European population structure. However, we lack a comprehensive assessment of historical genetic structure, including characterizing genetic heterogeneity and interactions across regions.

Several recent studies have begun sampling historical period genomes, including from Rome \citep{Antonio2019,Posth2021}, Iceland \citep{Ebenesersdottir2018}, and various European sites \citep{Fernandes2020,Olalde2019a}. These studies have revealed surprisingly high genetic diversity in cosmopolitan centers such as Imperial Rome. Integrating historical period genetics across a broad geographic range is instrumental to better understanding the development of European and Mediterranean population structure from prehistoric to present-day.

Here, we report 204 newly sequenced ancient genomes from across Europe and the Mediterranean, spanning the Iron Age through Late Antiquity. We characterize local population structure, identify ancestry outliers indicating long-distance mobility, and demonstrate that overall spatial population structure has remained remarkably stable over the last 3,000 years despite substantial individual-level mobility. Using spatial simulations, we show that the observed levels of long-range dispersal would be expected to erode population structure under standard models, suggesting that much of the observed mobility was transient in nature.

\section{Results}

\subsection{204 New Historical Genomes from Europe and the Mediterranean}

We collected whole genomes from 204 individuals across 53 archaeological sites in 18 countries spanning Europe and the Mediterranean (Figure~\ref{fig:timeline}), 26 of which were recently reported \citep{Moots2022}. These include the first published historical period genomes (after 1000 BCE) from present-day Armenia, Algeria, Austria, and France.

Dates for 126 samples were directly determined through radiocarbon dating, and these were used alongside archaeological contexts to infer dates for the remaining samples. DNA was extracted from either the powdered cochlear portion of the petrous bone ($n = 203$) or from teeth ($n = 1$). We performed whole-genome sequencing to a median depth of 0.92x (0.16x to 2.38x).

For downstream integration with published data, pseudohaploid genotypes were called for the 1240k SNP panel \citep{Mathieson2015}, resulting in a median of 685,058 SNPs (167,000 to 1,029,345) per sample. We analyzed newly reported genomes in conjunction with 2,033 present-day genomes, 1,998 prehistoric genomes, and 764 published historical period genomes \citep{Clemente2021,Mallick2020,Saupe2021}.

\subsection{Local Historical Population Structure Varies Across Regions}

We categorized the data into regional groups and characterized inter-individual heterogeneity by examining (1) variation of projections onto a PCA space of present-day genomes, (2) genetic groups identified by qpAdm and clustering across time within a region, and (3) admixture modeling of identified clusters \citep{Haak2015,Harney2021}.

PCA was performed on 829 individuals (480,712 SNPs) using smartpca v1600, with least-squares projection for ancient samples. A majority of regions have highly heterogeneous populations during the historical period. On average, we identified 10 genetic clusters within each region, with a minimum of 2 and a maximum of 23.

In Armenia, the population is highly homogeneous at any given time (Figure~\ref{fig:armenia}). After the Copper Age, two distinct genetic clusters are separated by a temporal split around 772--403 BCE. The earlier cluster (C1) includes newly reported samples ($n=5$) from Beniamin and published ones ($n=6$) from five other sites. The later cluster (C3), containing newly reported samples ($n=12$) from Beniamin dating between 403 BCE--500 CE, is genetically similar to present-day Armenians. Despite the split, there is evidence of partial continuity: the later cluster can be modeled using around 50\% of the earlier cluster plus an additional source of Steppe ancestry.

In contrast, newly collected samples reinforce previous findings of high heterogeneity in Rome \citep{Antonio2019,Posth2021}, including a large portion of the population having affinities for present-day Near Eastern populations. Southeastern European and Western European individuals during the Imperial Roman and Late Antiquity period also exhibit high heterogeneity, on par with that of contemporaneous Italy (Figures~\ref{fig:southeast} and~\ref{fig:western}).

In Southeastern Europe, a core group of individuals has ancestry similar to that of present-day and contemporaneous Central Europeans, while other clusters have ancestry similar to Northern Europeans and Eastern Mediterraneans. We also observe individuals of eastern nomadic ancestry, similar to previously reported Sarmatians, in both Western Europe ($n=2$) and Southeastern Europe ($n=2$).

\subsection{Ancestry Outliers and Cross-Regional Mobility}

We considered an individual an outlier if they belonged to an ancestry cluster that is underrepresented (fewer than 5\% of individuals or at most two individuals) within their sampling region from the Bronze Age to the present. We further identified outlier individuals who could be modeled as 100\% of a majority ancestry cluster found in a different region.

In total, we identified 11\% of individuals as outliers, and could connect 7\% of individuals to a putative source in a different region (Figure~\ref{fig:outliers}A). Based on the regions where these outliers and their sources originated, we created a network to illustrate their movements (Figure~\ref{fig:outliers}B).

No outliers were found within Armenia; however, we found outlier individuals in the Levant and Italy who can be putatively traced back to Armenia. In North Africa, outliers found in Iron Age Tunisia \citep{Moots2022} indicate movements from many regions in Europe, and North African-like outliers were found in Italy and Austria. One individual from Wels, Austria, located on the frontier of the Roman Empire, can be modeled using Canary Islander or Iron Age Tunisian individuals.

The 7\% estimate for outliers with source should be considered a lower bound. There are several cases where a cluster comprises more than 5\% of individuals in the region but are clearly of different ancestry than the majority and appear transient. No significant sex bias was detected among outliers.

\subsection{Spatial Population Structure is Relatively Stable in the Last 3,000 Years}

In each time period, we recovered the classical pattern of isolation-by-distance (Figure~\ref{fig:structure}A), where individuals closer in geographic space are also more similar genetically. Across time periods, there is a large decrease in overall $F_{ST}$ from the Mesolithic and Neolithic to the Bronze Age (approximately 10,000--2300 BCE), coinciding with the major prehistoric migrations \citep{Haak2015,Lazaridis2014}. From the Bronze Age onward, however, $F_{ST}$ does not decrease further with time, indicating that genetic differentiation across space is relatively stable from the Bronze Age to present-day.

Principal component analysis on 829 present-day European and Mediterranean genomes reveals close correspondence between genetic structure and geographic space, echoing previous observations \citep{Novembre2008}. We recovered similar spatial structure for historical samples, although noisier due to narrower sampling distribution and higher local genetic heterogeneity (Figure~\ref{fig:structure}B). The similarity in structure between present-day and historical period is especially striking in comparison to prehistoric genomes, which show much weaker correspondence to either present-day PCA or geographic space.

\subsection{Simulations Reveal That Observed Dispersal Would Collapse Structure}

Under standard population genetics models, is it surprising for stable population structure to be maintained in the presence of 7--11\% long-range migration? We addressed this question through spatial simulations, calibrating local dispersal to approximately match the observed $F_{ST}$ (maximum $\sim$0.03) (Figure~\ref{fig:simulations}A).

We then allowed a proportion of the population to disperse longer distances, empirically matching the migration distances observed in the data during the historical period. Even with long-range dispersal as low as 4\%, we observe decreasing $F_{ST}$ over 120 generations ($\sim$3,000 years). At 8\%, $F_{ST}$ decreases dramatically within 120 generations as spatial structure collapses to the point that it is hardly detectable in the first two principal components (Figure~\ref{fig:simulations}B--C).

These simulations indicate that under a basic spatial population genetics model, we would expect structure to collapse by the present day given the levels of movement observed. The persistence of population structure thus suggests a lower effective migration rate than indicated by observed dispersal.

\section{Discussion}

Our comprehensive analysis of 204 newly sequenced historical genomes, integrated with thousands of previously published ancient and modern genomes, reveals a paradox at the heart of European population history: remarkable individual-level mobility coexisting with remarkably stable population structure over the last 3,000 years.

The high heterogeneity observed within regions during the historical period reflects the extensively connected world of the Roman Empire and subsequent periods. Cross-Mediterranean contacts, military movements, trade networks, and forced migration created a cosmopolitan landscape where individuals of diverse genetic backgrounds lived side by side. Yet this did not translate into a collapse of the continental-scale population structure that was established during the Bronze Age.

We hypothesize that this paradox is explained by extensive transient dispersal. The Roman Empire's transportation networks and mobilization of people for trade, labor, and military service meant that many individuals moved long distances but did not necessarily establish permanent lineages in their new locations. Soldiers returned home, traders moved between ports, and enslaved individuals may not have had equivalent reproductive success. The result was a high observed rate of individual movement but a low effective migration rate in population genetics terms.

This work highlights the utility of ancient DNA in elucidating finer-scale human population dynamics in recent history, and demonstrates that even well-studied historical periods hold surprises when viewed through the lens of population genetics.

\section{Materials and Methods}

\subsection{Sample Collection}

Sampling was performed to maximize coverage across Europe and the Mediterranean, with particular focus on regions lacking published genomic data for the Imperial Roman and Late Antiquity period (e.g., France, Austria, the Balkans, Armenia, and North Africa). DNA was extracted from the powdered cochlear portion of the petrous bone or from teeth.

\subsection{Date Determination}

Dates were determined through radiocarbon dating for 126 samples, supplemented by archaeological context. Time periods were defined as: Mesolithic and Neolithic (10000--3500 BCE), Copper Age (3500--2300 BCE), Bronze Age (2300--1000 BCE), Iron Age (1000 BCE -- 1 CE), Imperial Roman and Late Antiquity (1 CE -- 700 CE), Medieval (700--1500 CE), and Post-Medieval/Present (1500 CE -- present).

\subsection{Population Genetic Analyses}

PCA was performed on 829 present-day individuals using smartpca v1600 with 5 outlier iterations and 10 principal components for outlier removal. Ancient genomes were projected onto this PCA space using least-squares projection. Ancestry modeling was performed using qpAdm. $F_{ST}$ was calculated between regional groups with confidence intervals from 200 bootstrap replicates.

\subsection{Spatial Simulations}

Simulations were performed to test whether observed levels of long-range dispersal would erode population structure. A base dispersal rate was calibrated to generate a maximum $F_{ST}$ of $\sim$0.03 across the maximal spatial distance using $N = 50{,}000$ and $\sigma_{Disp} = 0.02$. Long-range dispersal was added at rates of 4\% and 8\% with $\sigma_{DispLR} = 0.20$, and the effect on $F_{ST}$ was tracked over 120 generations.

\bibliography{references}

\end{document}
